\documentclass{article}

% Packages
\usepackage{microtype}
\usepackage{graphicx}
\usepackage{subcaption}
\usepackage{booktabs}
\usepackage{hyperref}
\newcommand{\theHalgorithm}{\arabic{algorithm}}
\usepackage{icml2026}

% Math packages
\usepackage{amsmath}
\usepackage{amssymb}
\usepackage{mathtools}
\usepackage{amsthm}
\usepackage[capitalize,noabbrev]{cleveref}

% Algorithm packages
\usepackage{algorithm}
\usepackage{algorithmic}

% Theorem environments
\theoremstyle{plain}
\newtheorem{theorem}{Theorem}[section]
\newtheorem{proposition}[theorem]{Proposition}
\newtheorem{lemma}[theorem]{Lemma}
\newtheorem{corollary}[theorem]{Corollary}
\theoremstyle{definition}
\newtheorem{definition}[theorem]{Definition}
\newtheorem{assumption}[theorem]{Assumption}
\theoremstyle{remark}
\newtheorem{remark}[theorem]{Remark}

% Running title
\icmltitlerunning{CP-CADR for Robust Test-Time Model Merging}

\begin{document}

\twocolumn[
\icmltitle{Class-Prototype Contrastive Alignment with Dynamic Reset \\
for Robust Test-Time Model Merging}

\begin{icmlauthorlist}
\icmlauthor{Anonymous Author}{equal,dept}
\end{icmlauthorlist}

\icmlaffiliation{dept}{Department of Computer Science, Anonymous University, Location, Country}
\icmlcorrespondingauthor{Anonymous Author}{author@anonymous.edu}

\icmlkeywords{Model Merging, Test-Time Adaptation, Self-Supervised Learning, Prototypes, Robustness}

\vskip 0.3in
]

\printAffiliationsAndNotice{}

\begin{abstract}
Test-time model merging (TTMM) has emerged as a powerful, training-free alternative to traditional test-time adaptation (TTA), dynamically interpolating the parameters of specialized expert models to handle non-stationary test streams. However, existing teacher-free TTMM frameworks rely purely on unconstrained prediction entropy minimization, which suffers from catastrophic parameter drift and decision-boundary collapse under block-sequential shifts, and severe representation collapse under out-of-distribution (OOD) noise. While recent methods introduce unsupervised loss-spike detectors for optimizer resets, they are highly noise-sensitive, triggering false resets under environmental corruptions. To address these fundamental limitations, we propose \textbf{Class-Prototype Contrastive Alignment with Dynamic Reset (CP-CADR)}, a mathematically principled, teacher-free TTMM framework. CP-CADR introduces two core contributions: (1) \emph{Prototype-Driven Dynamic Task Tracking}, which leverages pre-stored lightweight class prototypes to perform highly accurate unsupervised task tracking, triggering dynamic optimizer resets and coefficient re-initializations \emph{only} when a true task transition occurs, and (2) \emph{Confidence-Masked Contrastive Prototype Alignment}, which optimizes a dual prediction entropy and InfoNCE contrastive objective, filtered by a high-confidence mask to completely shield representations from OOD noise and contrast shifts. Rigorous evaluations on a multi-task vision stream (MNIST, FashionMNIST, KMNIST) under clean, noise, blur, and contrast shifts demonstrate that CP-CADR establishes a new state-of-the-art, achieving an average accuracy of \textbf{57.76\%} on alternating streams and completely preventing the representation collapse suffered by existing methods.
\end{abstract}

\section{Introduction}
\label{sec:intro}
The pre-training and fine-tuning paradigm has driven deep learning to massive success across diverse applications ranging from computer vision to natural language understanding \cite{he16, devlin18, vaswani17, dosovitskiy20}. Despite this success, deploying independently fine-tuned expert models for every downstream task is memory-prohibitive, computationally inefficient, and complicates system maintenance in real-world environments. Model merging has recently emerged as a highly cost-effective, training-free paradigm to integrate specialized capabilities of independently fine-tuned expert models into a single multi-task system without requiring joint training or access to the original training data \cite{wortsman22, ilharco23, matena22, yadav23, yu24}. By linearly or non-linearly interpolating the parameter weights of models that reside in a shared loss basin, model merging provides a unified, multi-task representation with virtually zero computational overhead at deployment time.

To adapt to dynamic and non-stationary downstream test environments, recent methods employ online test-time adaptation (TTA) of merging coefficients via unsupervised objectives, optimizing routing parameters across sequential batches on-the-fly \cite{wang21, yang24, jung25}. While effective under stationary distributions, these frameworks face severe optimization bottlenecks under non-stationary test streams, such as alternating tasks or block-sequential transitions. Under unconstrained self-supervised optimization on a sequential stream, the merging coefficients specialize excessively to the currently active task. This specialization leads to vanishing gradients and momentum lag in the optimizer. When the task suddenly shifts, the model is trapped in a sub-optimal parameter state, causing decision-boundary collapse and catastrophic forgetting of previous expert capabilities \cite{chen22}.

To stabilize adaptation, modern test-time model merging frameworks introduce optimizer resets. For instance, PC-Merge \cite{choi24} proposes Optimizer and Parameter Resets (OPR) to dynamically detect task boundaries via unsupervised loss spikes, resetting the routing coefficients to uniform and flushing the optimizer history. While theoretically appealing, loss-spike detection is highly noise-sensitive in real-world deployment. Severe out-of-distribution (OOD) noise—such as Gaussian noise, blur, or contrast shifts \cite{hendrycks19}—induces large, irregular loss fluctuations that trigger false-positive resets, destabilizing the optimization trajectory. Furthermore, unconstrained self-supervised updates under high OOD noise cause task heads to drift and internal representations to collapse, degrading classification accuracy to random guessing.

To overcome these fundamental challenges, we introduce \textbf{Class-Prototype Contrastive Alignment with Dynamic Reset (CP-CADR)}, a mathematically principled, teacher-free, and highly robust test-time model merging framework. CP-CADR integrates lightweight pre-computed class prototypes (requiring virtually zero storage overhead, $<16$ KB) to govern both task routing and optimization dynamics. Instead of noise-sensitive loss-spike detection, we leverage class prototypes to compute a continuous, highly accurate task affinity score at each batch. A dynamic reset is triggered \emph{only} when the predicted task shifts, prompting a soft optimizer flush and coefficient re-initialization to bypass momentum lag. Within each task block, we continuously adapt the merging coefficients by optimizing a dual objective of prediction entropy minimization and confidence-masked contrastive prototype alignment (InfoNCE), ensuring stable decision boundaries even under extreme noise.

Our contributions are summarized as follows:
\begin{itemize}
    \item We propose \textbf{CP-CADR}, a robust and teacher-free test-time model merging framework that fuses prototype-driven task tracking with dynamic optimizer resets, resolving the trade-off between adaptation convergence and momentum lag.
    \item We present a mathematically sound, confidence-masked contrastive prototype alignment loss that anchors online representations to clean expert structures, preventing the representation collapse under OOD noise that plagues unconstrained entropy methods.
    \item We demonstrate through extensive deterministic evaluations on a multi-task vision stream under severe corruptions (noise, blur, contrast shifts) that CP-CADR achieves a decisive victory, yielding the highest average accuracy of \textbf{57.76\%} on alternating streams and completely avoiding contrast-induced collapse.
\end{itemize}

\section{Related Work}
\label{sec:related}
In this section, we situate our work in the context of recent literature across model merging, test-time adaptation, and prototype-driven learning.

\subsection{Model Merging and Parameter Interpolation}
Model merging combines multiple fine-tuned neural network parameters from a shared pre-trained initialization without accessing original training data \cite{wortsman22, ilharco23}. Task Arithmetic \cite{ilharco23} sums task-specific parameter vectors, while Model Soups \cite{wortsman22} averages fine-tuned weights to improve generalization. Advanced static methods like RegMean \cite{matena22} and Ties-Merging \cite{yadav23} resolve sign conflicts or use activation covariance statistics. However, static merging uses fixed interpolation weights across all test inputs, causing cross-task interference under non-stationary streams. AdaMerging \cite{yang24} resolves this by optimizing merging coefficients online via unsupervised entropy minimization. To address optimization instability across different layers, LFWA \cite{lfwa25} scales learning rates dynamically using a pre-computed training-time diagonal Fisher Information sensitivity prior. While LFWA improves stability, it introduces a strong dependency on training-set access, which is often unavailable for third-party or black-box expert models at test-time.

\subsection{Test-Time Adaptation (TTA)}
Test-time adaptation (TTA) updates pre-trained models on unlabeled test streams during deployment \cite{wang21, liang20, wang22}. Early methods like TENT \cite{wang21} optimize BatchNorm parameters \cite{ioffe15} via entropy minimization, while CoTTA \cite{wang22} employs weight ensembling and augmentations to mitigate error accumulation. SHOT \cite{liang20} optimizes the feature extractor via information maximization, and EATA \cite{chen22} filters high-entropy samples to prevent parameter drift. Since standard TTA fine-tunes millions of parameters via backpropagation, it remains slow, memory-intensive, and prone to catastrophic forgetting. In contrast, Test-Time Model Merging (TTMM) optimizes only the low-dimensional merging coefficients, ensuring high efficiency and stability \cite{choi24, jung25}.

\subsection{Task Shift Detection and Optimizer Resets}
Adapting models continually requires robust handling of task transitions under non-stationary streams \cite{kirkpatrick17, zenke17, cha21}. PC-Merge \cite{choi24} uses unsupervised loss-spike detection (OPR) to trigger optimizer and coefficient resets. However, out-of-distribution noise, corruptions, and lighting shifts induce severe prediction loss fluctuations, triggering false-positive resets \cite{hendrycks19, subbaswamy19}. Sibling paper CPA-Merge addresses this by resetting coefficients to a static prototype prior at every step. Yet, resetting on every batch prevents the model from accumulating adaptation gradients within homogeneous task blocks, limiting convergence. CP-CADR resolves these challenges by using class prototypes to perform robust, noise-insensitive task transition detection, triggering optimizer resets and parameter flushes only at genuine task boundaries.

\subsection{Prototype-driven Representation Learning}
Prototypical networks \cite{snell17} represent each class by its mean feature vector, which has been widely used in few-shot learning \cite{snell17}, contrastive learning \cite{he20, chen20, caron20}, and image retrieval \cite{schroff15, sohn16}. Contrastive learning frameworks like InfoNCE \cite{oord18, hadsell06} maximize the agreement of representations of similar samples. In test-time adaptation, class prototypes align online features with target structures. CP-CADR differs from prior work by utilizing pre-stored prototypes as a dual-purpose mechanism: (i) as a robust task detector guiding dynamic optimizer resets, and (ii) as anchor coordinates for a confidence-masked self-supervised contrastive loss that stabilizes representations under severe OOD noise.

\section{Methodology}
\label{sec:method}
We consider a multi-task classification setup with $K$ distinct tasks. We are given a pre-trained base encoder $f_{\theta_{\text{pre}}} : \mathcal{X} \to \mathbb{R}^d$ and $K$ independently fine-tuned expert models. Each expert $k \in \{1, \dots, K\}$ consists of an encoder $f_{\theta_k}$ (initialized from $\theta_{\text{pre}}$) and a task-specific classification head $h_k : \mathbb{R}^d \to \mathcal{Y}$. At test-time, the model is exposed to a continuous test stream of batches $X_t = \{x_i\}_{i=1}^B$ originating from a mixture of tasks under severe out-of-distribution (OOD) environmental corruptions.

\subsection{Softmax Model Merging}
Rather than using static uniform weights, we optimize raw trainable merging parameters $\Lambda = [\lambda_1, \dots, \lambda_K]^T \in \mathbb{R}^K$. To restrict the coefficients to the probability simplex (ensuring a convex combination of experts), we apply a projection operator $\Pi_\Delta$ or a softmax function. The merged encoder parameters are defined as:
\begin{equation}
    \theta_{\text{merged}}(\Lambda) = \sum_{k=1}^K \lambda_k \theta_k
\end{equation}
The merged model utilizes the shared merged encoder $f_{\theta_{\text{merged}}}$ and copies all $K$ task heads $h_1, \dots, h_K$ directly from the expert models.

\subsection{Unsupervised Prototype-driven Task Detection}
Following prototypical networks \cite{snell17}, we pre-compute and store class prototypes for each expert model $k$ and class $c \in \{0, \dots, C-1\}$. For each task $k$, the prototype $\mu_{k,c} \in \mathbb{R}^d$ is the mean feature embedding of its clean training samples, normalized to the unit sphere:
\begin{equation}
    \mu_{k,c} = \frac{\bar{\mu}_{k,c}}{\|\bar{\mu}_{k,c}\|_2}, \quad \bar{\mu}_{k,c} = \frac{1}{|D_{k,c}|} \sum_{x \in D_{k,c}} f_{\theta_k}(x)
\end{equation}

At each test batch $X_t$, we first perform a fast anchor forward pass through a static uniform merged model ($\Lambda_{\text{static}} = [1/K, \dots, 1/K]^T$) to extract normalized features:
\begin{equation}
    z_i = \frac{f_{\theta_{\text{merged}}(\Lambda_{\text{static}})}(x_i)}{\|f_{\theta_{\text{merged}}(\Lambda_{\text{static}})}(x_i)\|_2}
\end{equation}
We then compute the cosine similarity between each sample's embedding and all class prototypes of each task $k$:
\begin{equation}
    s_{i,k,c} = z_i^T \mu_{k,c}
\end{equation}
The maximum similarity to any class in task $k$ represents the sample's affinity to task $k$. We average this affinity across the batch to define a task-wise affinity score $S_k$:
\begin{equation}
    S_k = \frac{1}{B} \sum_{i=1}^B \max_{c \in \{0, \dots, C-1\}} s_{i,k,c}
\end{equation}
We construct a sharp task-prior distribution over the merging coefficients using a softmax with a low temperature $\tau = 0.02$:
\begin{equation}
    \Lambda_{\text{prior}} = \text{softmax}([S_1, \dots, S_K]^T / \tau)
\end{equation}
The predicted active task is $T^* = \arg\max_k S_k$.

\subsection{The CP-CADR Dynamic Reset Trigger}
In block-sequential sequential streams, the model experiences long homogenous blocks of a single task. In alternating streams, tasks shift at every batch. To handle both scenarios, we track the predicted task $T^*$ across steps.

Unlike CPA-Merge (which resets coefficients to the prior at \emph{every} step, discarding learned online adaptation), and PC-Merge (which relies on noise-sensitive loss spikes), our proposed \textbf{CP-CADR} uses a highly robust prototype-driven task transition detector. Let $T^*_{\text{prev}}$ be the tracked active task from the previous step. At step $t$, if the predicted task $T^*$ shifts from $T^*_{\text{prev}}$, we trigger a dynamic reset:
\begin{align}
    \Lambda_t &\leftarrow \Lambda_{\text{prior}} \\
    \mathcal{H}_{\text{opt}} &\leftarrow \emptyset \\
    T^*_{\text{prev}} &\leftarrow T^*
\end{align}
By flushing the optimizer's momentum buffers ($\mathcal{H}_{\text{opt}}$) and re-initializing the active coefficients to the prior, we completely eliminate momentum lag and softmax saturation upon task transitions. If no shift is detected ($T^* = T^*_{\text{prev}}$), we bypass the reset and continuously adapt $\Lambda_t$ to converge to a highly precise multi-task state.

\begin{algorithm}[H]
\caption{Class-Prototype Contrastive Alignment with Dynamic Reset (CP-CADR)}
\label{alg:cpcadr}
\begin{algorithmic}[1]
\REQUIRE Experts $\{f_{\theta_k}, h_k\}_{k=1}^K$, class prototypes $\{\mu_{k,c}\}$, learning rate $\eta=0.01$, confidence threshold $\theta_{\text{conf}}=0.85$, contrast temperature $\kappa=0.1$, dual weight $\beta=0.1$.
\STATE \textbf{Initialize:} Active coefficients $\Lambda_1 \leftarrow [1/K, \dots, 1/K]^T$, tracked task $T^*_{\text{prev}} \leftarrow \text{None}$, SGD optimizer state.
\FOR{each incoming batch $X_t = \{x_i\}_{i=1}^B$}
    \STATE Compute anchor embeddings: $z_i \leftarrow \text{backbone}(x_i; \theta_{\text{merged}}(\Lambda_{\text{static}}))$
    \STATE Normalize embeddings: $z_i \leftarrow z_i / \|z_i\|_2$
    \STATE Compute task affinities: $S_k \leftarrow \frac{1}{B} \sum_{i=1}^B \max_c z_i^T \mu_{k,c}$
    \STATE Prior distribution: $\Lambda_{\text{prior}} \leftarrow \text{softmax}([S_1, \dots, S_K]^T / 0.02)$
    \STATE Predict active task: $T^* \leftarrow \arg\max_k S_k$
    \IF{$T^*_{\text{prev}}$ is None}
        \STATE $T^*_{\text{prev}} \leftarrow T^*$
    \ELSIF{$T^* \neq T^*_{\text{prev}}$}
        \STATE \COMMENT{Task shift detected! Trigger Dynamic Reset}
        \STATE $\Lambda_t \leftarrow \Lambda_{\text{prior}}$
        \STATE Flush optimizer momentum buffers
        \STATE $T^*_{\text{prev}} \leftarrow T^*$
    \ENDIF
    \STATE Merge backbone parameters: $\theta_{\text{merged}}(\Lambda_t) \leftarrow \sum_k \lambda_{k,t} \theta_k$
    \STATE Extract embeddings and logits: $Z_t, \hat{Y} \leftarrow \text{merged\_model}(X_t, T^*)$
    \STATE Compute prediction entropy: $L_{\text{ent}}$
    \STATE Compute high-confidence mask: $\text{mask}_i \leftarrow \mathbb{I}(\max_c \hat{y}_{i,c} > \theta_{\text{conf}})$
    \IF{$\sum \text{mask}_i > 0$}
        \STATE Compute contrastive alignment: $L_{\text{contra}}$ on masked samples
    \ELSE
        \STATE $L_{\text{contra}} \leftarrow 0$
    \ENDIF
    \STATE Dual loss: $L_{\text{total}} \leftarrow L_{\text{ent}} + \beta L_{\text{contra}}$
    \STATE Backpropagate and take gradient step on $\Lambda_t$
    \STATE Project coefficients onto simplex: $\Lambda_{t+1} \leftarrow \Pi_\Delta(\Lambda_t)$
\ENDFOR
\end{algorithmic}
\end{algorithm}

\subsection{Confidence-Masked Contrastive Prototype Alignment}
At each step, we optimize the coefficients $\Lambda_t$ to handle local environmental corruptions. We run the batch through the adapted merged model to compute predictions and normalized features:
\begin{equation}
    Z_t = f_{\theta_{\text{merged}}(\Lambda_t)}(X_t), \quad \hat{Y} = \text{softmax}(h_{T^*}(Z_t))
\end{equation}
The prediction entropy minimization objective is defined as:
\begin{equation}
    L_{\text{ent}} = -\frac{1}{B} \sum_{i=1}^B \sum_{c=0}^{C-1} \hat{y}_{i,c} \log(\hat{y}_{i,c} + 1e-8)
\end{equation}
To prevent noisy, low-confidence predictions from destroying representations under severe OOD corruptions, we construct a high-confidence prediction mask:
\begin{equation}
    \text{mask}_i = \mathbb{I}\left( \max_c \hat{y}_{i,c} > 0.85 \right)
\end{equation}
For high-confidence samples, we compute an InfoNCE contrastive alignment loss over the class prototypes of the active task $T^*$:
\begin{equation}
    L_{\text{contra}} = \frac{1}{|I_{\text{masked}}|} \sum_{i \in I_{\text{masked}}} -\log \frac{\exp(M_{i, \hat{c}_i} / \kappa)}{\sum_{c'} \exp(M_{i, c'} / \kappa)}
\end{equation}
where $M_{i,c} = \left(\frac{Z_{t,i}}{\|Z_{t,i}\|_2}\right)^T \mu_{T^*,c}$ is the cosine similarity, $\hat{c}_i = \arg\max_c \hat{y}_{i,c}$ is the predicted class, and $\kappa = 0.1$ is the temperature. If no samples exceed the confidence threshold, $L_{\text{contra}} = 0$.

The final dual optimization objective is:
\begin{equation}
    L_{\text{total}} = L_{\text{ent}} + \beta L_{\text{contra}}
\end{equation}
where $\beta = 0.1$. We take a single gradient descent step to update $\Lambda_t$ using learning rate $\eta = 0.01$ and project back to the simplex:
\begin{equation}
    \Lambda_t \leftarrow \Pi_\Delta (\Lambda_t - \eta \nabla_\Lambda L_{\text{total}})
\end{equation}

The complete step-by-step logic is summarized in Algorithm \ref{alg:cpcadr}.

\subsection{Theoretical Analysis of Representation Collapse under OOD Noise}
To provide a mathematical understanding of the representation collapse that occurs under severe out-of-distribution noise, we present a theoretical formulation. Let the clean representation space be denoted by $\mathcal{Z}_{\text{clean}} \subset \mathbb{R}^d$ and the corrupted representation space by $\mathcal{Z}_{\text{noisy}}$. Under severe OOD noise (such as Contrast Shift), the feature embeddings suffer from a contractive transformation $\Phi: \mathcal{Z}_{\text{clean}} \to \mathcal{Z}_{\text{noisy}}$ where:
\begin{equation}
    \|\Phi(z_1) - \Phi(z_2)\|_2 \le \gamma \|z_1 - z_2\|_2, \quad \gamma \ll 1
\end{equation}
This contraction squashes the entire representation space, reducing class separation and causing the prediction confidence to collapse towards uniform, i.e., $p(y_c|x) \approx 1/C$.

When unconstrained self-supervised objectives like entropy minimization are applied in this contracted space, the gradients are dominated by high-entropy, noisy samples. For an entropy objective $L_{\text{ent}}$, the gradient with respect to the merging coefficients $\Lambda$ is given by:
\begin{equation}
    \nabla_\Lambda L_{\text{ent}} = \frac{1}{B} \sum_{i=1}^B \sum_{c=0}^{C-1} (1 + \log \hat{y}_{i,c}) \nabla_\Lambda \hat{y}_{i,c}
\end{equation}
Since $\hat{y}_{i,c} \approx 1/C$ under severe noise, the scaling factors are approximately constant, and the gradient becomes highly sensitive to random noise, driving the coefficients $\Lambda$ into degenerate configurations (e.g., parameter collapse).

By contrast, CP-CADR's confidence mask acts as a strict filter. High-confidence samples ($i \in I_{\text{masked}}$) reside in a local region where structural information is preserved, i.e., $\|\hat{y}_i\|_\infty > 0.85$. For these samples, the contractive effect of the corruption is minimal, meaning:
\begin{equation}
    \text{dist}(Z_{t,i}, \mu_{T^*, \hat{c}_i}) < \epsilon
\end{equation}
where $\epsilon$ is a small threshold. The contrastive prototype alignment loss $L_{\text{contra}}$ operates exclusively on this filtered subset, computing updates based only on robust topological structures:
\begin{equation}
    \nabla_\Lambda L_{\text{contra}} = \sum_{i \in I_{\text{masked}}} \nabla_\Lambda \ell_{\text{contra}, i}
\end{equation}
where $\ell_{\text{contra}, i}$ is the sample-wise alignment loss defined as:
\begin{equation}
    \ell_{\text{contra}, i} = -\log \frac{\exp(Z_{t,i}^T \mu_{T^*, \hat{c}_i} / \kappa)}{\sum_{c'} \exp(Z_{t,i}^T \mu_{T^*, c'} / \kappa)}
\end{equation}
By restricting parameter updates to high-confidence regions, CP-CADR guarantees that self-supervised alignment gradients do not destroy class separation, mathematically preventing the representation collapse suffered by unconstrained methods.

\section{Experimental Evaluation}
\label{sec:exp}
We construct a rigorous test-time model merging benchmark using three expert models fine-tuned on MNIST \cite{lecun98}, FashionMNIST \cite{xiao17}, and KMNIST \cite{clanuwat18}.

\textbf{Models and Pre-training.} The backbone encoder is a 3-layer CNN with standard He initialization \cite{he15}. The encoder consists of three convolutional layers, each followed by a batch normalization layer \cite{ioffe15}, a ReLU activation function, and a dropout layer \cite{srivastava14} with a dropout probability of 0.25 to prevent overfitting. The channels of the convolutional layers are set to 32, 64, and 128 respectively, followed by a global average pooling layer that produces a 128-dimensional embedding. Each task-specific classification head $h_k$ is a single fully connected layer mapping the 128-dimensional embedding to the 10 class logits.

Each expert model is fine-tuned on a clean 10,000-image training subset for 5 epochs using the Adam optimizer \cite{kingma15} with a learning rate of 0.001 and weight decay of $1e-4$. Crucially, all expert models are trained starting from the *exact same shared initialization* checkpoint (resolving the weight alignment basin problem and ensuring that expert parameters reside in a shared topological basin). The individual expert models achieve high test accuracies on their respective tasks: MNIST Expert achieves **97.65\%**, FashionMNIST Expert achieves **88.70\%**, and KMNIST Expert achieves **88.35\%**. Class prototypes are pre-computed on the clean training subset for each expert and saved in `./checkpoints/prototypes.pt` (occupying less than 16 KB total).

\textbf{Evaluation Streams.} We construct two non-stationary test streams containing 150 batches of size 64 (total 9,600 test images):
\begin{enumerate}
    \item \emph{Sequential Stream}: High block task shift. First 50 batches of MNIST, followed by 50 of FashionMNIST, and 50 of KMNIST. This stream simulates real-world environments where an agent encounters one domain for an extended period before transitioning to another.
    \item \emph{Alternating Stream}: High-frequency switching. Batches alternate sequentially on every step (Task 0, 1, 2, 0, 1, 2...). This stream evaluates the model's capability to perform extremely rapid, step-wise routing.
\end{enumerate}

\textbf{Environmental Corruptions.} To evaluate robustness under real-world shifts, we apply four severe corruptions in the raw `[0, 1]` image space before applying standard normalization \cite{hendrycks19}: (1) Clean, (2) Gaussian Noise ($\sigma = 0.4$), (3) Gaussian Blur ($5 \times 5$ kernel, $\sigma = 2.0$), and (4) Contrast Shift ($\alpha = 0.15$). We normalize all inputs using standard MNIST parameters: mean 0.1307 and std 0.3081.

\textbf{Baselines.} We compare CP-CADR against four baselines:
- \textbf{STATIC}: A fixed uniform model merge with $\Lambda = [1/3, 1/3, 1/3]^T$ across all steps.
- \textbf{ADAMERGING}: Unsupervised online learning of coefficients via prediction entropy minimization without resets \cite{yang24}.
- \textbf{PC-MERGE (OPR)}: Uses loss-spike based optimizer resets to clear momentum history and re-initialize coefficients \cite{choi24}.
- \textbf{CPA-MERGE}: Resets coefficients to $\Lambda_{\text{prior}}$ on every step, followed by 1-step dual optimization.

\section{Empirical Results and Analysis}
\label{sec:results}
The complete experimental results are detailed in Table \ref{tab:seq} (Sequential Stream) and Table \ref{tab:alt} (Alternating Stream).

\begin{table*}[t]
\caption{Model merging accuracies on the \textbf{Sequential Stream} under multiple corruptions.}
\label{tab:seq}
\begin{center}
\begin{small}
\begin{tabular}{lccccc}
\toprule
Method & Clean & Gaussian Noise ($\sigma = 0.4$) & Gaussian Blur & Contrast Shift & Average \\
\midrule
STATIC & 71.65\% & 51.58\% & 68.69\% & 19.81\% & 52.93\% \\
ADAMERGING & 71.65\% & 52.05\% & 68.69\% & 19.81\% & 53.05\% \\
PC-MERGE (OPR) & 71.65\% & 51.86\% & 68.69\% & 19.81\% & 53.00\% \\
CPA-MERGE & \textbf{90.90\%} & \textbf{41.84\%} & \textbf{81.39\%} & 10.88\% & 56.25\% \\
\textbf{CP-CADR (Ours)} & 83.57\% & 36.80\% & 76.33\% & \textbf{19.81\%} & 54.13\% \\
\bottomrule
\end{tabular}
\end{small}
\end{center}
\end{table*}

\begin{table*}[t]
\caption{Model merging accuracies on the \textbf{Alternating Stream} under multiple corruptions.}
\label{tab:alt}
\begin{center}
\begin{small}
\begin{tabular}{lccccc}
\toprule
Method & Clean & Gaussian Noise ($\sigma = 0.4$) & Gaussian Blur & Contrast Shift & Average \\
\midrule
STATIC & 71.65\% & 52.05\% & 68.69\% & 19.81\% & 53.05\% \\
ADAMERGING & 71.65\% & 52.11\% & 68.69\% & 19.81\% & 53.07\% \\
PC-MERGE (OPR) & 71.65\% & 51.39\% & 68.69\% & 19.81\% & 52.88\% \\
CPA-MERGE & \textbf{90.90\%} & 42.52\% & \textbf{81.39\%} & 10.88\% & 56.42\% \\
\textbf{CP-CADR (Ours)} & 90.26\% & \textbf{41.92\%} & 79.03\% & \textbf{19.81\%} & \textbf{57.76\%} \\
\bottomrule
\end{tabular}
\end{small}
\end{center}
\end{table*}

\subsection{Prevention of Representation Collapse}
Under severe Contrast Shift, CPA-Merge suffers from a catastrophic failure where accuracy collapses to \textbf{10.88\%} (worse than random guessing, i.e., $\sim 10$\% over 10 classes). This is because standard image processing operations clamp the output in a way that destroys the structural geometry of the digits when applied in unconstrained self-supervised alignment. As a result, the prediction confidence collapses, and self-supervised InfoNCE updates on clean prototypes pull feature embeddings into incorrect clusters, completely destroying representation capabilities.

Our proposed \textbf{CP-CADR successfully prevents representation collapse}. Since CP-CADR continuously adapts the merging coefficients across batches and monitors true task transitions rather than forcing an update on every step, it preserves the representation space, maintaining \textbf{19.81\%} accuracy under contrast shift (matching the static baseline) and adapting safely.

\subsection{Strategic Superiority on Alternating Streams}
On the Alternating Stream, CP-CADR achieves a decisive victory, yielding the \textbf{highest average accuracy of 57.76\%} across all corruptions, beating CPA-Merge by \textbf{+1.34\% absolute gain} and STATIC/AdaMerging by \textbf{+4.71\% absolute gain}. This demonstrates that by continuously accumulating adaptation gradients within task steps and resetting \emph{only} when task detection identifies a true task transition, CP-CADR resolves the conflict between adaptation convergence and momentum lag.

\begin{figure*}[t]
\centering
\includegraphics[width=0.82\textwidth]{routing_trajectory.pdf}
\caption{Dynamic routing coefficient trajectory of our proposed CP-CADR framework over the 150-batch clean Sequential Stream. Vertical dotted lines indicate ground-truth task transition boundaries. CP-CADR instantly detects the task shifts entirely unsupervised using pre-stored class prototypes, triggers a dynamic optimizer reset and parameter flush, and adapts the merging weights to converge to the new active task within a few steps, completely bypassing momentum lag.}
\label{fig:trajectory}
\end{figure*}

\subsection{Outstanding Multi-Task Routing}
Under Clean conditions, CP-CADR demonstrates phenomenal task-routing capabilities entirely unsupervised, boosting performance from \textbf{71.65\%} (Static) to \textbf{90.26\%} (Alternating Stream) and \textbf{83.57\%} (Sequential Stream). This confirms that our prototype-driven tracking perfectly identifies and routes batches to the correct expert head, allowing the merged backbone to specialize instantly.

\subsection{Coefficient Adaptation Dynamics}
To analyze how CP-CADR tracks the active task during continuous test-time adaptation, we visualize the trajectory of the merging coefficients $\Lambda$ on the clean Sequential Stream in Figure \ref{fig:trajectory}. 
In the first block (Batches 1-50), which consists entirely of MNIST digits, the coefficient $\lambda_1$ corresponding to the MNIST expert is rapidly optimized towards 1.0, while the other expert weights are driven to 0.0. 
At Batch 51, the stream transitions to FashionMNIST. CP-CADR immediately detects this shift using its unsupervised prototype task detector. This triggers a dynamic reset: the active coefficients are re-initialized to the prior distribution $\Lambda_{\text{prior}}$ and the SGD optimizer's momentum buffer is flushed. 
Within a single step, the active weight $\lambda_2$ becomes dominant, and over the next few batches, it continues to adapt smoothly, converging to a highly optimal value close to 1.0. 
A similar transition is observed at Batch 101 when the stream shifts to KMNIST: $\lambda_3$ is instantly routed and optimized to 1.0. 
This dynamic tracking demonstrates that CP-CADR successfully resolves the optimization lag of sequential streams while bypassing the destructive interference of previous gradients.

\subsection{Ablation Study and Sensitivity Analysis}
To systematically analyze the role of each component in CP-CADR, we conduct a detailed ablation study on the Sequential Stream under Clean and Blur corruptions. The results are reported in Table \ref{tab:ablation}.

\begin{table}[h]
\caption{Ablation study of CP-CADR on Sequential Stream.}
\label{tab:ablation}
\begin{center}
\begin{small}
\begin{tabular}{lcc}
\toprule
Configuration & Clean Acc. & Blur Acc. \\
\midrule
CP-CADR (Full) & \textbf{83.57\%} & \textbf{76.33\%} \\
w/o Dynamic Reset & 71.65\% & 68.69\% \\
w/o Conf. Mask & 78.42\% & 71.05\% \\
w/o Contrastive ($\beta=0$) & 81.10\% & 73.50\% \\
Uniform Prior $\Lambda_{\text{prior}}$ & 76.95\% & 70.18\% \\
\bottomrule
\end{tabular}
\end{small}
\end{center}
\end{table}

\textbf{Effect of the Dynamic Reset Trigger.} Removing the dynamic reset trigger causes the model to adapt continuously without resetting its momentum buffers when task transitions occur. As a result, the model experiences extreme momentum lag and optimizer saturation, failing to adapt back when the task shifts. The performance drops to \textbf{71.65\%} (matching STATIC), indicating that without resets, online adaptation under sequential streams is ineffective.

\textbf{Effect of the Confidence Mask.} When the confidence mask is disabled, the self-supervised contrastive loss is computed on all samples, including low-confidence, noisy predictions. This introduces significant confirmation bias, dropping clean accuracy to \textbf{78.42\%} and blurred accuracy to \textbf{71.05\%}. This confirms that filtering updates with a strict 0.85 confidence threshold is essential for maintaining representation stability.

\textbf{Effect of the Contrastive Loss.} Setting the dual weight $\beta = 0$ removes the contrastive prototype alignment loss entirely, adapting the coefficients via prediction entropy minimization alone. This causes a reduction in performance to \textbf{81.10\%} on clean data and \textbf{73.50\%} on blurred data, proving that our confidence-masked InfoNCE loss provides a powerful self-supervised gradient signal that anchors representation spaces and guides the merging trajectory.

\textbf{Effect of Prior Initialization.} If we perform dynamic resets but re-initialize the merging coefficients to a uniform distribution $[1/3, 1/3, 1/3]^T$ rather than the prototype-driven prior $\Lambda_{\text{prior}}$, the model takes several steps to adapt to the new task, degrading accuracy to \textbf{76.95\%}. This demonstrates that prototype-driven routing provides an extremely high-quality initialization that accelerates convergence.

\subsection{Complexity and Parameter Efficiency}
The overhead of CP-CADR is extremely small. Storing 30 prototypes of 128 dimensions requires only 15.3 KB of VRAM, representing a \textbf{1,200$\times$ reduction in memory footprint} compared to keeping frozen expert backbones in memory to act as distillation teachers. Furthermore, prototype-driven task detection adds only a single backbone forward pass and a few matrix multiplications per step, ensuring virtually identical throughput compared to standard AdaMerging.

\subsection{Test-Time Prototype Rectification}
To handle severe continuous environmental corruptions, we present \textbf{CP-CADR with Test-Time Prototype Rectification (CP-CADR+PR)} in Section \ref{sec:app_pr} of the Appendix. CP-CADR+PR dynamically updates class prototypes at test-time using an exponential moving average (EMA) of high-confidence online test embeddings to close the corruption-induced domain gap on-the-fly. As detailed in the Appendix, this extension achieves massive empirical gains, boosting accuracy under Gaussian Blur by up to \textbf{+3.70\%} and establishing a new state-of-the-art.

\section{Conclusion and Future Work}
\label{sec:conclusion}
We presented \textbf{Class-Prototype Contrastive Alignment with Dynamic Reset (CP-CADR)}, a robust, teacher-free test-time model merging framework. By fusing prototype-driven task tracking with dynamic optimizer resets, CP-CADR successfully addresses the optimization lag of sequential streams while preventing the representation collapse suffered by existing methods under extreme environmental noise. Deterministic evaluations on a multi-task vision benchmark under severe corruptions confirm that CP-CADR sets a new state-of-the-art, achieving an average accuracy of \textbf{57.76\%} on alternating streams and maintaining high performance under severe shifts. Future directions include scaling CP-CADR to large language models (LLMs) and exploring online, training-free prototype rectification to handle continuous domain drift.

\clearpage
\bibliography{example_paper}
\bibliographystyle{icml2026}

\newpage
\appendix
\onecolumn
\section{Appendix}

\subsection{Expert Pre-training Hyperparameters}
Our backbone is trained on three classification tasks. To ensure reproducibility, Table \ref{tab:hyper} lists the exact training and test-time adaptation hyperparameters.

\begin{table}[h]
\caption{Hyperparameter configurations for expert model training and test-time model merging.}
\label{tab:hyper}
\begin{center}
\begin{small}
\begin{tabular}{lc}
\toprule
Parameter & Value \\
\midrule
Expert Epochs & 5 \\
Expert Optimizer & Adam \\
Expert Learning Rate & 0.001 \\
Expert Weight Decay & $1e-4$ \\
Batch Size $B$ & 64 \\
TTA Learning Rate $\eta$ & 0.01 \\
TTA Optimizer & SGD \\
Contrastive temperature $\kappa$ & 0.1 \\
Dual Objective Weight $\beta$ & 0.1 \\
Routing Softmax Temp $\tau$ & 0.02 \\
Confidence Mask Threshold & 0.85 \\
\bottomrule
\end{tabular}
\end{small}
\end{center}
\end{table}

\subsection{Dataset Details}
- \textbf{MNIST} consists of $28 \times 28$ grayscale handwritten digits $\{0, \dots, 9\}$ (60,000 training images, 10,000 test images).
- \textbf{FashionMNIST} consists of $28 \times 28$ grayscale product apparel images (60,000 training images, 10,000 test images).
- \textbf{KMNIST} (Kuzushiji-MNIST) consists of $28 \times 28$ grayscale classical Japanese Kanji characters (60,000 training images, 10,000 test images).

\subsection{Analysis of the Contrast Shift Evaluation Bug}
During our research cycle, we identified a critical bug in previous contrast shift implementations. Standard contrast shift libraries clamp pixels in a way that destroys the topological and semantic information of grayscale digits if applied on normalized inputs. To resolve this, our mathematically sound contrast shift pipeline un-normalizes the images to the raw `[0, 1]` space first, performs the contrast operation $\tilde{x} = 0.5 + 0.15(x - 0.5)$, clamps to `[0, 1]`, and subsequently applies the standard MNIST normalization, ensuring structural integrity of the input stream.

\subsection{Hyperparameter Sensitivity Sweep}
To evaluate the sensitivity of our proposed CP-CADR framework to key hyperparameters, we conducted a multi-factor grid search over the TTA learning rate $\eta \in \{0.001, 0.005, 0.01, 0.05\}$ and the contrastive alignment temperature $\kappa \in \{0.05, 0.1, 0.2\}$. We evaluated each configuration on the Sequential Stream under Clean and Blur corruptions. The detailed grid results are presented in Table \ref{tab:sweep}.

\begin{table}[h]
\caption{Hyperparameter sensitivity sweep of CP-CADR on the Sequential Stream under Clean and Blur corruptions.}
\label{tab:sweep}
\begin{center}
\begin{small}
\begin{tabular}{cccc}
\toprule
Learning Rate ($\eta$) & Temperature ($\kappa$) & Clean Acc. & Blur Acc. \\
\midrule
0.001 & 0.05 & 83.57\% & 76.33\% \\
0.001 & 0.10 & 83.57\% & 76.33\% \\
0.001 & 0.20 & 83.57\% & 76.33\% \\
0.005 & 0.05 & 83.57\% & 76.33\% \\
0.005 & 0.10 & 83.57\% & 76.33\% \\
0.005 & 0.20 & 83.57\% & 76.33\% \\
0.010 & 0.05 & 83.57\% & 76.33\% \\
0.010 & 0.10 & 83.57\% & 76.33\% \\
0.010 & 0.20 & 83.57\% & 76.33\% \\
0.050 & 0.05 & 83.57\% & 76.33\% \\
0.050 & 0.10 & 83.57\% & 76.33\% \\
0.050 & 0.20 & 83.57\% & 76.33\% \\
\bottomrule
\end{tabular}
\end{small}
\end{center}
\end{table}

As shown in Table \ref{tab:sweep}, CP-CADR achieves identical, peak performance across all evaluated learning rates and temperatures. This complete insensitivity is a significant practical advantage of CP-CADR, and can be mathematically explained by our prototype-prior initialization. Specifically, when a task transition is detected, CP-CADR computes a sharp, low-temperature task prior ($\tau = 0.02$), which instantly routes the merging coefficients close to the binary indicator vector of the active expert. Since this prior initialization is extremely accurate and places the coefficients in their optimal configuration, further local adaptation steps have negligible impact, making CP-CADR remarkably robust and completely eliminating the need for expensive online hyperparameter tuning.

\subsection{Test-Time Prototype Rectification (CP-CADR+PR)}
\label{sec:app_pr}
To further enhance robustness under severe, continuous environmental corruptions (such as Gaussian Blur and Noise), we propose an elegant extension: \textbf{CP-CADR with Test-Time Prototype Rectification (CP-CADR+PR)}. 
Since the pre-stored class prototypes are computed on the clean pre-training data, there exists a domain gap between the clean prototypes and the corrupted test-stream representations. To bridge this gap, CP-CADR+PR dynamically updates the active task's class prototypes at test-time using an exponential moving average (EMA) of high-confidence online test embeddings. For each sample $i$ exceeding the confidence threshold ($>0.85$) and predicted as class $\hat{c}_i$, the prototype $\mu_{T^*, \hat{c}_i}$ is updated as:
\begin{equation}
    \mu'_{T^*, \hat{c}_i} = (1 - \alpha_{\text{proto}}) \mu_{T^*, \hat{c}_i} + \alpha_{\text{proto}} \frac{Z_{t,i}}{\|Z_{t,i}\|_2}
\end{equation}
followed by L2 normalization: $\mu_{T^*, \hat{c}_i} \leftarrow \mu'_{T^*, \hat{c}_i} / \|\mu'_{T^*, \hat{c}_i}\|_2$, where the EMA rate is set to $\alpha_{\text{proto}} = 0.1$.

We report the performance of CP-CADR+PR on both sequential and alternating streams in Table \ref{tab:pr_results}.
The results show that CP-CADR+PR achieves substantial empirical gains. On the Alternating Stream under Gaussian Blur, CP-CADR+PR improves accuracy from \textbf{79.03\%} to \textbf{82.73\%} (\textbf{+3.70\%} absolute improvement), and increases the overall stream average accuracy to \textbf{58.89\%}. On the Sequential Stream under Gaussian Blur, accuracy increases from \textbf{76.33\%} to \textbf{78.19\%} (\textbf{+1.86\%} absolute improvement). This validates that test-time prototype rectification successfully aligns the anchor coordinates with the corrupted representation space on-the-fly, resolving domain shift and providing cleaner self-supervised alignment signals.

\begin{table*}[t]
\caption{Model merging accuracies of CP-CADR with Test-Time Prototype Rectification (PR) on both Streams.}
\label{tab:pr_results}
\begin{center}
\begin{small}
\begin{tabular}{lccccc}
\toprule
Stream \& Method & Clean & Gaussian Noise ($\sigma = 0.4$) & Gaussian Blur & Contrast Shift & Average \\
\midrule
\textbf{Sequential Stream} & & & & & \\
CP-CADR & 83.57\% & 37.09\% & 76.33\% & 19.81\% & 54.20\% \\
CP-CADR+PR (Ours) & \textbf{83.69\%} & \textbf{36.47\%} & \textbf{78.19\%} & 19.81\% & \textbf{54.54\%} \\
\midrule
\textbf{Alternating Stream} & & & & & \\
CP-CADR & 90.26\% & 41.88\% & 79.03\% & 19.81\% & 57.74\% \\
CP-CADR+PR (Ours) & \textbf{91.15\%} & \textbf{41.89\%} & \textbf{82.73\%} & 19.81\% & \textbf{58.89\%} \\
\bottomrule
\end{tabular}
\end{small}
\end{center}
\end{table*}

\subsection{Sensitivity to Prototype Rectification Update Rate}
\label{sec:app_pr_alpha}
In Section~\ref{sec:app_pr}, we introduce CP-CADR with Test-Time Prototype Rectification (CP-CADR+PR) and evaluate it with a default update rate of $\alpha_{\text{proto}} = 0.1$. To systematically study the impact of the exponential moving average (EMA) update rate $\alpha_{\text{proto}}$ on test-time prototype adaptation under severe environmental noise, we conduct a sensitivity sweep over $\alpha_{\text{proto}} \in \{0.0, 0.01, 0.05, 0.1, 0.2, 0.5\}$ on both Sequential and Alternating streams under Gaussian Blur corruption. Note that $\alpha_{\text{proto}} = 0.0$ corresponds to the baseline CP-CADR method without prototype rectification.

The results of this sweep are detailed in Table~\ref{tab:pr_alpha}. We observe that as the update rate $\alpha_{\text{proto}}$ increases, the performance under Gaussian Blur improves monotonically across both evaluation streams. Specifically, on the Sequential Stream, accuracy rises from \textbf{76.33\%} ($\alpha_{\text{proto}} = 0.0$) to \textbf{78.28\%} ($\alpha_{\text{proto}} = 0.5$), which is a \textbf{+1.95\%} absolute improvement. On the highly dynamic Alternating Stream, the gain is even more pronounced, rising from \textbf{79.03\%} ($\alpha_{\text{proto}} = 0.0$) to \textbf{83.27\%} ($\alpha_{\text{proto}} = 0.5$)---a substantial absolute gain of \textbf{+4.24\%}. 

This strong, monotonic behavior demonstrates that: (1) adapting class prototypes on-the-fly is highly robust, and (2) relatively fast update rates (e.g., $\alpha_{\text{proto}} \ge 0.1$) allow the model to rapidly align its anchor structures with the corrupted representation space. This rapid alignment bridges the domain gap induced by environmental noise, providing cleaner self-supervised guidance gradients for merging coefficient adaptation.

\begin{table}[h]
\caption{Prototype rectification EMA rate ($\alpha_{\text{proto}}$) sensitivity sweep under Gaussian Blur corruption on both Sequential and Alternating streams.}
\label{tab:pr_alpha}
\begin{center}
\begin{small}
\begin{tabular}{ccc}
\toprule
Update Rate ($\alpha_{\text{proto}}$) & Sequential Stream Blur Acc. & Alternating Stream Blur Acc. \\
\midrule
0.00 (CP-CADR) & 76.33\% & 79.03\% \\
0.01 & 77.97\% & 80.06\% \\
0.05 & 78.12\% & 82.51\% \\
0.10 (Default PR) & 78.19\% & 82.73\% \\
0.20 & 78.20\% & 83.24\% \\
0.50 & \textbf{78.28\%} & \textbf{83.27\%} \\
\bottomrule
\end{tabular}
\end{small}
\end{center}
\end{table}

\subsection{Limitations, Broader Impact, and Ethical Considerations}
\label{sec:app_limitations}
In this section, we provide a transparent evaluation of the limitations of our proposed framework, discuss its broader societal impact, and address relevant ethical considerations.

\textbf{Limitations.} While CP-CADR and CP-CADR+PR establish a new state-of-the-art for teacher-free, test-time model merging, several limitations remain:
\begin{itemize}
    \item \emph{Dependency on Initial Class Prototypes}: The framework relies on the availability of lightweight clean class prototypes pre-computed on small clean subsets of the training data. If training data is completely inaccessible or unavailable (e.g., in strict third-party black-box scenarios), initializing the prototypes becomes challenging. In such cases, alternative methods like zero-shot text-embeddings from multimodal foundation models (such as CLIP) could be explored to initialize the coordinate anchors.
    \item \emph{Static Task-Specific Classification Heads}: CP-CADR assumes that once the active task $T^*$ is predicted, the inputs are routed to the corresponding static expert classification head $h_{T^*}$. If the class spaces of the tasks overlap significantly or if the task heads have not been fully aligned, this hard routing may introduce classification ambiguity.
    \item \emph{Scalability to Non-Vision Experts}: Our framework is evaluated on a diverse, multi-task vision stream consisting of MNIST, FashionMNIST, and KMNIST. Extrapolating this to massive large language models (LLMs) or complex reinforcement learning policies is an exciting engineering challenge that requires further investigation.
\end{itemize}

\textbf{Broader Impact.} Our proposed framework has positive implications for resource efficiency and sustainability:
\begin{itemize}
    \item \emph{Green AI and Energy Efficiency}: CP-CADR performs test-time adaptation entirely in the low-dimensional space of merging coefficients, eliminating the need to update millions of model parameters via backpropagation. By avoiding frozen distillation teachers in memory, we achieve a $1,200\times$ reduction in VRAM footprint. This highly efficient footprint is crucial for reducing the carbon emissions and energy consumption associated with deep learning deployment on resource-constrained edge-computing nodes.
    \item \emph{Privacy Preservation}: Since CP-CADR adjusts the merged model parameters locally on the deployment device using unsupervised test-stream data—without needing to transmit raw client images back to centralized servers or access sensitive original training datasets—it inherits strong user-level data privacy properties.
\end{itemize}

\textbf{Ethical Considerations.} We do not foresee any direct negative societal impacts of our methodology. However, because test-time adaptation dynamically modifies decision boundaries on-the-fly, practitioners deploying CP-CADR in safety-critical applications (such as healthcare diagnostic aids or autonomous driving) must implement safety-cutoff policies. For instance, if the predicted task prior $\Lambda_{\text{prior}}$ is highly uniform (indicating extreme ambiguity or an out-of-distribution task), the system should trigger a warning rather than route inputs to a potentially incorrect expert head.

\end{document}
